\documentclass{article}
% Double-blind workshop submission. NOT [final]: that is camera-ready and
% de-anonymises the author block.
\usepackage[dblblindworkshop]{neurips_2026}
\workshoptitle{Interpretability for Discovery}

\usepackage[utf8]{inputenc}
\usepackage[T1]{fontenc}
\usepackage{hyperref}
\usepackage{url}
\usepackage{booktabs}
\usepackage{amsmath}
\usepackage{graphicx}
\usepackage{microtype}
\usepackage{array}

\title{Two Readouts of a Traffic GNN, and Why Neither Supports a Claim About the
Road Network}

\author{%
  Anonymous Author(s)\\
  Affiliation\\
  Address\\
  \texttt{email}
}

\begin{document}
\maketitle

\begin{abstract}
We ask whether two standard post-hoc readouts of a spatio-temporal GNN --- a
mask-based explainer and the model's own learned adjacency --- yield checkable
structure about the road network they were trained on. Neither does. What makes
that null worth reporting rather than merely asserting is the instrumentation:
three independent stability checks, an external geometric baseline that is
independent of the model, a second explainer family, and a pre-registered
reliability proxy that failed its own criteria. The same geometric test that finds
nothing for the explainer recovers road structure at $0.76$--$0.81$ precision for a
distance baseline, so the instrument demonstrably detects real structure when it is
present. We also report one apparent effect that a sample-size robustness check did
not reproduce, and withdraw it.
\end{abstract}

\section{Introduction}

Interpretability becomes discovery only when a readout says something about the
world that can be checked independently of the model that produced it. We take a
spatio-temporal GNN trained to forecast traffic speed on the METR-LA sensor network
and ask directly: does reading its internals tell us anything testable about the
road network? Two readouts are available and both are standard --- a GNNExplainer-style
soft mask over input sensors, and the adaptive adjacency the model learns for itself
--- and we find that neither supports such a claim. The result is a null, so the
instrumentation carries the weight: we measure stability three independent ways,
compare against an \emph{external} check that does not involve the model at all
(physical road distance between sensors, taken from the road-distance table the
dataset ships), cross-check against a second explainer family, and pre-register a
cheap reliability proxy that then fails. That external baseline is what makes the
null interpretable rather than uninformative: applied to a nearest-by-road rule it
recovers the road graph at high precision, so a near-zero score for the explainer
reflects the explainer, not a weak test.

\section{Setup}

The forecaster is a Graph WaveNet-style ST-GNN~\citep{gwnet2019} on METR-LA (207
loop sensors). Its spatial support mixes a \emph{given} adjacency with a
\emph{learned} one as $A_{\text{final}} = \alpha A_{\text{phys}} + (1-\alpha)
A_{\text{sem}}$. The given adjacency is a thresholded Gaussian kernel on directed
road distance: with $\sigma = 2584.5$\,m and threshold $\kappa = 0.1$ the effective
cutoff is $3921.7$\,m, giving 1{,}515 directed edges over 207 nodes (density
$0.0355$). ``Off-adjacency'' below therefore means precisely ``more than about
3.9\,km apart by road,'' not a vague notion of distance. The learned component is
$A_{\text{sem}} = \mathrm{softmax}(\mathrm{relu}(EE^\top))$ with $E$ a learned
$[207, 10]$ embedding initialised from $\mathcal{N}(0, 0.01^2)$ --- not seeded from
the kernel --- and on the checkpoint we analyse the model puts weight $0.7032$ on the
learned graph against $0.2968$ on the given one. The first readout is a
GNNExplainer-style~\citep{ying2019gnnexplainer} node mask, optimised per (target,
window) and read as a top-$k$ ($k{=}8$) source list; the aggregate over windows is
an influence matrix $W$. Unless stated otherwise every number below comes from one
run of $207 \times 24 = 4{,}968$ mask solves on CPU, and from the same epoch-54
checkpoint, so the two readouts describe one model.\footnote{Explainer solves: run
\texttt{20260824T235016Z\_\_influence\_\allowbreak graph\_solves}. Learned adjacency:
run \texttt{20260825T093132Z\_\_learned\_\allowbreak semantic\_graph}, checkpoint
\texttt{metr\_la\_best.pt}, sha \texttt{046b2a4a}, epoch 54.}

\section{Result 1: the mask is close to flat}

Covering $80\%$ of a target's off-self importance mass takes on average
$146.6 \pm 4.7$ sources in the congested regime and $154.0 \pm 2.9$ in free flow,
out of 206 available. A readout that needs roughly three quarters of the network to
account for four fifths of its mass is not selecting a small causal set, and the
narrow standard deviations show this is a property of the method rather than of a
few odd targets. The learned adjacency is no better: its rows need $142.1$ of 206
sources for the same $80\%$, it is structurally dense (every entry positive), and
each row spreads over 107 of 207 nodes --- $52\%$ of uniform. The model's spatial
receptive field is effectively the whole graph in one hop.

\section{Result 2: instability, measured three ways}

Flat readouts should be unstable, and they are. Table~\ref{tab:stability} collects
the evidence under perturbations that have nothing to do with one another.

\begin{table}[h]
\centering
\caption{Stability of the two readouts. Every entry is agreement between two
computations that should have given the same answer. ``Pooled'' is one Jaccard over
the union of all targets' top-8 source--target pairs across the full node set, with
no adjacency or reachability restriction; ``per-target mean'' averages each target's
own Jaccard and is a \emph{different} statistic, hence the spread. Chance references
are two independent uniform draws under the same construction.}
\label{tab:stability}
\small
\setlength{\tabcolsep}{4pt}
\begin{tabular}{@{}llcc@{}}
\toprule
Perturbation & Statistic & Value & Chance ref. \\
\midrule
\multicolumn{4}{l}{\emph{Mask explainer}}\\
Window resample, free-flow (200 targets) & pooled top-8 Jaccard & 0.2256 & 0.0165 \\
Window resample, congested (28 targets)  & pooled top-8 Jaccard & 0.1399 & 0.0182 \\
Window resample, $\lambda/3\lambda/10\lambda$ (40) & per-target mean & 0.2323 / 0.2385 / 0.2287 & --- \\
Optimiser seed, $\lambda/3\lambda/10\lambda$ (10) & per-target mean & 0.3480 / 0.3149 / 0.3015 & --- \\
CPU vs MPS, same seed (3 cases) & top-8 Jaccard & 0.3991 & --- \\
CPU vs MPS, same seed (3 cases) & Pearson, magnitudes & 0.9193 & --- \\
\midrule
\multicolumn{4}{l}{\emph{Learned adjacency}}\\
Epoch 34 vs 54, same run & top-8 Jaccard & 0.2163 & 0.0163 \\
Against a separate training run & top-8 Jaccard & 0.1532 / 0.1589 & 0.0163 \\
\midrule
\multicolumn{4}{l}{\emph{Between the two readouts}}\\
$W$ vs $A_{\text{sem}}$, free-flow / congested & per-target top-8 J & 0.0873 / 0.1743 & --- \\
\bottomrule
\end{tabular}
\end{table}

Three things are worth drawing out. First, raising the sparsity penalty tenfold does
not concentrate the mask --- per-target agreement is flat across $\lambda$, at
$0.2323 \pm 0.1461$, $0.2385 \pm 0.1739$ and $0.2287 \pm 0.1831$ --- so the
instability is not a tuning artefact that a stronger prior would remove. Second, the
CPU-versus-MPS row is the starkest: the same target, the same window and the same
seed agree on importance \emph{magnitudes} at Pearson $0.9193$ (Spearman $0.5778$)
while agreeing on top-8 \emph{identity} only $0.3991$ of the time ($0.3576$ at
top-20). Floating-point associativity alone reorders which sensors a reader would be
told mattered. Third, the learned adjacency preserves \emph{ranks} across training
states far better than it preserves top-$k$ identity (per-target Spearman $0.7795$
between epochs against a top-8 Jaccard of $0.2163$), which is the same dissociation
seen in the mask: what is stable is not what gets reported. The two readouts also
largely disagree with each other, with 84 of 200 free-flow targets sharing
\emph{nothing} in their top 8 (17 of 28 when congested) and per-target Spearman of
$+0.0618$ and $-0.115$.

\section{Result 3: the external geometric check}

The checkable question is whether either readout recovers road geometry. We score
each method's top-8 edges against the given kernel adjacency on the same targets
(Table~\ref{tab:geometry}).

\begin{table}[h]
\centering
\caption{Precision against the kernel adjacency, same targets and same $k{=}8$ for
every row. The nearest-by-road baseline is the calibration: it must score well, since
the adjacency is a road-distance threshold, and the fact that it does is what shows
the test detects road structure when a method encodes it.}
\label{tab:geometry}
\small
\begin{tabular}{lcc}
\toprule
Method & Free-flow & Congested \\
\midrule
Nearest by road distance (baseline) & 0.7557 & 0.8088 \\
Mask explainer, top-8                & 0.2863 & 0.1397 \\
Degree-matched random                & 0.0344 & 0.0551 \\
Uniform random                       & 0.0331 & 0.0368 \\
\bottomrule
\end{tabular}
\end{table}

The explainer sits above both random floors, so its mask is not noise --- but it
recovers a small fraction of what a plain distance rule recovers on the identical
test. \textbf{That gap is the paper's calibration point}: the instrument is
demonstrably able to detect road structure, so a low score for the explainer reflects
the explainer rather than a weak test. The learned adjacency does worse still. Its
top-8 overlaps a nearest-by-road rule at $0.0644$, and it places $0.0445$ of its mass
on adjacency edges against a $0.0355$ chance rate --- a lift of $1.252$, close to
none. Nor does the explainer favour long-range structure, which would have been a
different result worth reporting: beyond eight chained kernel radii it is
\emph{depleted}, at $0.570\times$ the base rate in free flow and $0.923\times$ when
congested.

\section{Result 4: a second explainer family disagrees, invisibly}

On a 28-scenario subset where SHAP (KernelExplainer) was run against the same model,
the two methods' top-$k$ sets overlap at a mean Jaccard of $0.0221$. Two principled
methods for reading one model agree on about $2\%$ of what they nominate. Yet a
downstream consumer cannot tell: a deterministic citation checker scoring an LLM's
narration of each returns \emph{identical} precision $0.9821$ and hallucination
$0.0179$ for both.\footnote{\texttt{shap\_comparison\_summary.json}; $n{=}28$.} Whatever
disagreement exists between explainer families is not visible to anything built on
top of them. \emph{Checkpoint caveat}: this comparison ran on the epoch-34 checkpoint,
not the epoch-54 one used elsewhere in this paper. Both of its arms use
that same checkpoint, so the $0.0221$ overlap is an internally consistent
within-model comparison; but it is a different training state from Results 1--3 and
we do not pool it with them. The summary file records no checkpoint, so we
established this by re-solving cached explanations on both candidates: they
reproduce exactly on epoch 34 (top-8 Jaccard $1.000$) and not on epoch 54
($0.067$ and $0.000$).

\section{Result 5: an apparent effect, withdrawn}

We report one result that did not survive its own robustness check, because the
check is the point.

Restricting to explainer edges beyond eight chained kernel radii and asking which
survive both split halves, 52 of 423 free-flow edges survive. Those survivors sit in
the learned adjacency's own top-8 at a rate of $0.1538$, against $0.0404$ for edges
surviving only one half --- an apparent shift of roughly four-fold, which would have
been the most interesting thing here: two independent readouts agreeing about
long-range structure that the road graph does not contain.

Recomputing on the first 64 targets of the same seeded shuffle --- an unbiased
prefix, so this isolates sample size rather than geography --- the shift did not reproduce: survivors sit in the top-8 at $0.0714$ against $0.0381$ for
one-half-only, on 14 survivors instead of 52. The congested regime shows $0.0$ at
both sample sizes. We therefore withdraw the effect. Two readings are consistent
with the data --- the prefix is underpowered at 14 survivors, or the full-sample
shift is not robust --- and neither is established, so the honest statement is that
the effect has not been shown to replicate and the four-fold figure should not be
quoted on its own. We flag the direction explicitly because it is the opposite of
the more familiar failure: here the apparent effect is at the \emph{larger} sample
and the subsample is what fails to reproduce it, which makes it ambiguous rather
than cleanly negative. The next section is the cleanly negative case.

\section{Result 6: a pre-registered reliability proxy, falsified}

If flatness drives instability then mask entropy should predict it, and entropy costs
one solve where measuring stability costs two. We fixed the criteria before running:
Spearman $\rho < -0.30$ against the split-half Jaccard \emph{and} a bootstrap CI
excluding zero, plus a threshold chosen on the sweep alone flagging all 12 committed
explanations. Pooled over 120 (target, setting) pairs, $\rho = -0.2008$ with 95\% CI
$[-0.4229, +0.0452]$ for the entropy of the window-averaged mask, and
$\rho = -0.1732$, CI $[-0.3938, +0.0513]$ for the mean per-window entropy --- the
latter being the only form obtainable from a single solve, and so the only one that
would have made the proxy cheap. Both intervals include zero and the threshold
flagged 3 of 11 distinct committed explanations. Both criteria are falsified. A
single-setting version at $n{=}40$ did appear to pass, at $\rho = -0.3412$ with 95\%
CI $[-0.5999, -0.0241]$, meeting both criteria; it did not survive the full
pool.\footnote{Recomputed into its own run directory
\texttt{20260827T110720Z\_\_entropy\_\allowbreak gate\_single\_setting} from the
Stage-1 masks, using the same functions as the pooled analysis, because the original
inline value had been overwritten by the pooled result.} That is the expected fate of an
underpowered positive, and the reason the criteria were fixed in advance.

\section{Discussion}

Three things would make either readout usable for a claim about the road network:
stability under perturbations that should not matter, agreement across explainer
families, and agreement with a check external to the model. None holds here. Top-$k$
identity moves under window resampling, optimiser seed, and floating-point
associativity alike; two explainer families overlap at $2\%$; and the external
geometric test, calibrated by a baseline that passes it comfortably, finds little.

The general lesson is about consumption, not about this model. An interpretability
output that is handed downstream --- to an LLM that narrates it, to an automatic
checker that scores the narration, to a human reading a ranked list --- carries its
instability with it, and nothing downstream surfaces that. The SHAP comparison is the
sharpest version: two readouts that disagree almost completely produce byte-identical
downstream scores. Any pipeline that turns an attribution into a claim should
therefore carry its own stability estimate alongside the attribution, or the claim
inherits an error bar nobody has measured.

\section{Responsible use statement}

This work concerns interpretability readouts of a traffic-forecasting model, in a
research setting, on a public sensor dataset with no personal data. Our results
should not be read as a safety claim about the model or the explainer for their
original purpose: we did not evaluate forecasting accuracy here, and the absence of a
validated discovery does not imply the forecaster is unfit for forecasting or that
the explainer is unfit for the debugging and sanity-checking uses it is normally put
to. The scoped claim is narrower and concerns downstream use: these particular
readouts should not be treated as evidence about real road structure --- for instance
to justify a traffic-management intervention at a sensor named in a top-$k$ list ---
without the stability and external checks demonstrated here. Presenting an unstable
attribution to an operator as though it were a determined fact is the specific harm
this paper argues against.

\section{Limitations}

\textbf{(a)} One dataset. Everything here is METR-LA; we make no cross-city claim,
and whether these nulls hold on another sensor network is untested. \textbf{(b)} The
external baseline inherits a definition. ``Geometry'' here means the dataset's
directed road-distance table thresholded at $3921.7$\,m, which is one sensible notion
and not the only one; moreover only $27.1\%$ of ordered sensor pairs have a road
distance recorded at all ($24.4\%$ of non-adjacent pairs), so the baseline is
evaluated on the covered subset. \textbf{(c)} Some checks are small: the seed-repeat
uses 10 targets, the cross-device comparison 3 cases, the SHAP comparison 28
scenarios, and the congested split-half only 28 targets. These are indicative. The
central claims rest on the full-sample split-half (200 free-flow targets) and the
geometry comparison, both at $n{=}207$ solved targets. \textbf{(d)} A single
explainer family plus one SHAP cross-check; other attribution methods might behave
differently, which is exactly what Result 4 suggests is worth testing.

\section{Related work}

Evaluating GNN explainers against ground truth is itself contested.
\citet{sanchez2020evaluating} benchmark attribution against computable ground-truth
attributions on synthetic graphs and \citet{agarwal2023evaluating} supply generators
and a library for the same purpose, while \citet{faber2021ground} show that comparing
explanations to ground truth is unsound when the trained GNN does not use the
ground-truth edges, separating the responsible, causing and explanatory motifs. Our
setting has no ground-truth attribution at all, which is why we substitute an
external \emph{geometric} check and calibrate it with a baseline that passes.

The second readout we examine is the adaptive adjacency now standard in
spatio-temporal forecasting: Graph WaveNet~\citep{gwnet2019} introduced a learned
self-adaptive matrix, MTGNN~\citep{wu2020connecting} learns a unidirectional graph
from node embeddings, and AGCRN~\citep{bai2020adaptive} learns node-specific
parameters and infers inter-series dependencies without a predefined graph. These are
often described as capturing hidden spatial dependency, which invites reading the
learned matrix as a discovered graph. Our measurements argue against that reading for
this model: the learned matrix is structurally dense, spreads $52\%$ of uniform
across each row, aligns with road adjacency at a lift of $1.252$, and its top-8
changes substantially between training states.

\begin{thebibliography}{9}
\bibitem[Agarwal et al.(2023)]{agarwal2023evaluating}
C.~Agarwal, O.~Queen, H.~Lakkaraju, and M.~Zitnik.
\newblock Evaluating explainability for graph neural networks.
\newblock \emph{Scientific Data}, 10:144, 2023.

\bibitem[Bai et al.(2020)]{bai2020adaptive}
L.~Bai, L.~Yao, C.~Li, X.~Wang, and C.~Wang.
\newblock Adaptive graph convolutional recurrent network for traffic forecasting.
\newblock In \emph{NeurIPS}, 2020.

\bibitem[Faber et al.(2021)]{faber2021ground}
L.~Faber, A.~K.~Moghaddam, and R.~Wattenhofer.
\newblock When comparing to ground truth is wrong: On evaluating GNN explanation
methods.
\newblock In \emph{KDD}, 2021.

\bibitem[Wu et al.(2019)]{gwnet2019}
Z.~Wu, S.~Pan, G.~Long, J.~Jiang, and C.~Zhang.
\newblock Graph WaveNet for deep spatial-temporal graph modeling.
\newblock In \emph{IJCAI}, pages 1907--1913, 2019.

\bibitem[Wu et al.(2020)]{wu2020connecting}
Z.~Wu, S.~Pan, G.~Long, J.~Jiang, X.~Chang, and C.~Zhang.
\newblock Connecting the dots: Multivariate time series forecasting with graph
neural networks.
\newblock In \emph{KDD}, 2020.

\bibitem[Sanchez-Lengeling et al.(2020)]{sanchez2020evaluating}
B.~Sanchez-Lengeling, J.~N.~Wei, B.~K.~Lee, E.~Reif, P.~Wang, W.~W.~Qian,
K.~McCloskey, L.~J.~Colwell, and A.~B.~Wiltschko.
\newblock Evaluating attribution for graph neural networks.
\newblock In \emph{NeurIPS}, 2020.

\bibitem[Ying et al.(2019)]{ying2019gnnexplainer}
R.~Ying, D.~Bourgeois, J.~You, M.~Zitnik, and J.~Leskovec.
\newblock GNNExplainer: Generating explanations for graph neural networks.
\newblock In \emph{NeurIPS}, 2019.
\end{thebibliography}

\appendix

\section{Claim-to-source map}
\label{app:sources}

Every number in the main text traces to one of the following. Run directories are
under \texttt{evaluation/results/}; \texttt{IG} abbreviates
\texttt{20260824T235016Z\_\_influence\_graph\_solves}, \texttt{LG} abbreviates
\texttt{20260825T093132Z\_\_learned\_semantic\_graph}, and \texttt{GWI} abbreviates
\texttt{20260825T201433Z\_\_grounding\_without\_information}.

\footnotesize
\begin{tabular}{@{}>{\raggedright\arraybackslash}p{0.34\linewidth}>{\raggedright\arraybackslash}p{0.24\linewidth}>{\raggedright\arraybackslash}p{0.36\linewidth}@{}}
\toprule
Claim & Value & Source \\
\midrule
Kernel adjacency construction & $\sigma$ 2584.5\,m, $\kappa$ 0.1, cutoff 3921.7\,m, 1515 edges, density 0.0355 & IG \texttt{metrics.json}, \texttt{geometry} \\
$\alpha$ split, learned vs given & 0.7032 / 0.2968 & LG \texttt{learned\_graph.json}, \texttt{per\_checkpoint.\allowbreak metr\_la\_best.pt.alpha} \\
Mask flatness, 80\% mass & 146.6$\pm$4.7 congested; 154.0$\pm$2.9 free-flow, of 206 & IG \texttt{metrics.json}, \texttt{per\_regime.*.\allowbreak sources\_for\_mass\_frac} \\
Learned graph flatness & 142.1 of 206; 107 of 207 per row (52\% of uniform) & LG, \texttt{sources\_for\_mass\_frac}, \texttt{density} \\
Split-half J, \emph{pooled} & 0.2256 (200 targets) / 0.1399 (28); ref 0.0165 / 0.0182 & IG \texttt{metrics.json}, \texttt{per\_regime.*.\allowbreak stability.jaccard\_all\_edges} \\
Split-half J, \emph{per-target mean}, 3 $\lambda$ & 0.2323 / 0.2385 / 0.2287 & GWI \texttt{part4\_metrics.json}, \texttt{settings[]} \\
Seed-repeat J, 3 $\lambda$ & 0.3480 / 0.3149 / 0.3015 ($n{=}10$) & GWI \texttt{part4\_metrics.json} \\
CPU vs MPS & Pearson 0.9193; top-8 J 0.3991; Spearman 0.5778; top-20 J 0.3576 & \texttt{20260824T235643Z\_\_\allowbreak device\_divergence}, \texttt{device\_divergence.json} \\
Cross-checkpoint J, learned graph & 0.2163 / 0.1532 / 0.1589; ref 0.0163 & LG, \texttt{cross\_checkpoint\_stability} \\
$W$ vs $A_{\text{sem}}$ agreement & J 0.0873 (84 zero) / 0.1743 (17 zero); Spearman +0.0618 / $-$0.115 & IG \texttt{metrics.json}, \texttt{per\_regime.*.\allowbreak vs\_learned\_semantic\_graph} \\
Geometry precision, explainer & 0.2863 free-flow / 0.1397 congested & IG \texttt{metrics.json}, \texttt{per\_regime.*.edges.top\_k.\allowbreak precision} \\
Geometry precision, baselines & nearest-road 0.7557 / 0.8088; uniform 0.0331 / 0.0368; degree 0.0344 / 0.0551 & IG \texttt{metrics.json}, \texttt{per\_regime.*.baselines} \\
Learned graph vs geometry & overlap 0.0644; mass 0.0445 vs 0.0355 chance, lift 1.252 & LG, \texttt{overlap\_with\_nearest\_road}, \texttt{road\_alignment\_semantic} \\
Beyond-8-radii enrichment & 0.57$\times$ free-flow; 0.923$\times$ congested & IG \texttt{metrics.json}, \texttt{per\_regime.*.beyond\_k} \\
SHAP overlap and downstream scores (\textbf{epoch 34}) & J 0.0221; precision 0.9821; halluc 0.0179; $n{=}28$ & \texttt{shap\_comparison\_summary.json}; epoch established by re-solve, see \S6 \\
Survivors, full sample & 52 of 423; 0.1538 vs 0.0404; survival 0.1229 & IG \texttt{metrics.json}, \texttt{per\_regime.free\_flow.\allowbreak beyond\_k\_survivors} \\
Survivors, 64-target prefix & 14 survivors; 0.0714 vs 0.0381; congested 0.0 both & IG \texttt{metrics.json}, \texttt{survivor\_sample\_\allowbreak size\_sensitivity} \\
Entropy proxy $\rho$ and CIs & $-$0.2008 [$-$0.4229, +0.0452]; $-$0.1732 [$-$0.3938, +0.0513]; 3/11 & GWI \texttt{verdicts.json}, \texttt{P4} \\
Entropy proxy at $n{=}40$ & $\rho$ $-$0.3412, CI [$-$0.5999, $-$0.0241] & \texttt{20260827T110720Z\_\_entropy\_\allowbreak gate\_single\_setting}, \texttt{entropy\_gate\_\allowbreak single\_setting.json} \\
Road-distance coverage & 27.1\% of pairs; 24.4\% of non-adjacent & IG \texttt{metrics.json}, \texttt{geometry} \\
\bottomrule
\end{tabular}
\normalsize

\end{document}
